% !TEX program = xelatex
% !TeX TXS-program:compile = txs:///xelatex/[-shell-escape]
\documentclass[x11names,svgnames,usepdftitle=false]{beamer}

\usetheme{metropolis}
\setbeamertemplate{caption}{\raggedright\insertcaption\par}
\setbeamercolor{block title}{fg=LightSkyBlue4}
\newenvironment<>{importantblock}[1]{%
  \begin{actionenv}#2%
    \def\insertblocktitle{#1}%
    \par%
    \mode<presentation>{%
      \setbeamercolor{block title}{%
        use=normal text,
        fg=normal text.fg,
        bg=normal text.bg!80!fg
      }
      \setbeamercolor{block body}{
        use={block title, normal text},
        bg=block title.bg!50!normal text.bg
      }
    }%
    \usebeamertemplate{block begin}}
  {\par\usebeamertemplate{block end}\end{actionenv}}
\newenvironment<>{decisionblock}[1][Statement of the problem]{%
  \begin{actionenv}#2%
    \def\insertblocktitle{#1}%
    \par%
    \mode<presentation>{%
      \setbeamercolor{block title}{%
        use=normal text,
        fg=normal text.fg,
        bg=normal text.bg!80!fg
      }
      \setbeamercolor{block body}{
        use={block title, normal text},
        bg=block title.bg!50!normal text.bg
      }
    }%
    \usebeamertemplate{block begin}}
  {\par\usebeamertemplate{block end}\end{actionenv}}
\usepackage{fontspec}
\usepackage[french, frenchkw]{algorithm2e}
\SetKwFor{Pc}{Pour chaque}{faire}{fin}
\SetKwProg{Fn}{Fonction}{:}{}
\usepackage{textcomp}
\usepackage{minted}
\usepackage[framemethod=TikZ]{mdframed}
\usepackage{pdfpages}
\usepackage{qrcode}
\usepackage[francais]{babel}
\providecommand{\og}{«\,}
\providecommand{\fg}{\,»}
\usepackage{xspace}
\usepackage[normalem]{ulem}
\usepackage[backend=biber,style=authoryear,url=false]{biblatex}
\addbibresource{Bibliography.bib} % Bibliography file
\usepackage{amsmath}
\usepackage{mathtools}
\usepackage{pifont}% http://ctan.org/pkg/pifont
\newcommand{\cmark}{\ding{51}}%
\newcommand{\xmark}{\ding{55}}%
\usepackage{marvosym}
\usepackage{tikz}
\usetikzlibrary{calc, positioning, arrows, automata, fadings, decorations.pathmorphing, shapes, patterns, tikzmark, fit, chains, snakes, shapes.multipart}
\setminted{fontsize=\small}
% ------------------------------------------------------------------
% Couleurs et puces des slides refaites en natif (morpion, itérables)
% ------------------------------------------------------------------
\definecolor{cmblue}{RGB}{86, 185, 238}
\definecolor{cmgold}{RGB}{244, 197, 66}
\definecolor{cmmint}{RGB}{224, 245, 240}
\colorlet{idxcol}{blue!60!black}% indice / position
\colorlet{eltcol}{red!70!black}% élément / valeur d'une case
% petit tableau 3 x 3 (lu ligne par ligne) : #1 = les 9 éléments séparés par des virgules, #2 = code TikZ ajouté (surlignages)
\newcommand{\tableau}[2]{%
  \begin{tikzpicture}[x=1cm,y=1cm,baseline=(current bounding box.north)]
    \draw[thick, rounded corners=2pt] (0,0) rectangle (3.0,-3.0);
    \foreach \t in {1,2}{\draw[thick] (\t,0) -- (\t,-3.0); \draw[thick] (0,-\t) -- (3.0,-\t);}
    \foreach \e [count=\k from 0] in {#1}{%
      \pgfmathtruncatemacro{\r}{int(\k/3)}\pgfmathtruncatemacro{\c}{mod(\k,3)}%
      \node[font=\ttfamily\Large, text=eltcol] at ({\c+0.5},{-\r-0.5}) {\e};}
    \foreach \i in {0,1,2}{%
      \node[font=\small\bfseries, text=idxcol] at ({\i+0.5},0.3) {\i};
      \node[font=\small\bfseries, text=idxcol] at (-0.3,{-\i-0.5}) {\i};}
    #2
  \end{tikzpicture}%
}
\newcommand{\hlcell}[2]{\draw[cmgold, very thick, rounded corners=3pt] ({#2+0.05},{-#1-0.05}) rectangle ({#2+0.95},{-#1-0.95});}
\newcommand{\hlrow}[1]{\draw[cmgold, very thick, rounded corners=3pt] (0.05,{-#1-0.05}) rectangle (2.95,{-#1-0.05-0.9});}
% puce : tiret bleu
\newcommand{\tiret}{\textcolor{cmblue}{\rule[0.35ex]{0.8em}{0.4ex}}\hspace{0.5em}}

% https://tex.stackexchange.com/questions/401884/how-do-i-change-hyperlinks-color-only
\hypersetup{
  colorlinks,
  allcolors=.,
  urlcolor=DarkBlue,
  pdftitle=Algorithmique et Programmation 1 - CM6
}


\author{}
\date{Lundi 19 octobre 2026}

\title{\includegraphics[width=0.75\textwidth]{img/logo-igm.png} \\ Algorithmique et Programmation 1}
\institute{L1 Mathématiques - L1 Informatique \\ Semestre 1}

\begin{document}

\maketitle

\begin{frame}{CC1 : Comment et quoi ?}
  \begin{block}{Horaire}
    \vspace{0pt}
    \textbf{Lundi 2 novembre 14h-16h} (tiers-temps : 16h40) en \alert{salles à déterminer}
  \end{block}
  \pause
  \begin{block}{Contenu}
    \vspace{-0.25em}
    \begin{itemize}
    \item Tout ce qu'on a vu, y compris les séances de cette semaine
    \item Les slides \alert{et} les notes de cours
    \item QCM
    \item Questions de programmation écrites
    \item Mini-problème
    \end{itemize}
  \end{block}
  \pause
  \begin{block}{Sur la triche}
    \vspace{-0.25em}
    \begin{itemize}
    \item L'examen est \textbf{individuel}
    \item Tolérance zéro
    \end{itemize}
  \end{block}
\end{frame}

\begin{frame}{Questions ouvertes et mini-problème}
  Voir \href{DOC-catalog.pdf}{le sujet de l'an dernier}
\end{frame}


\section{Approfondissements : les listes \newline (de listes (de listes (\dots)))}

% Slide refaite en natif (c'était l'image slide_morpion.png)
\begin{frame}{Revenons à notre morpion}
  \begin{tikzpicture}[x=1cm,y=1cm]
    \useasboundingbox (0,0) rectangle (10.6,6.2);
    \fill[cmblue] (0,0) rectangle (10.6,6.2);
    \node[text=white, font=\fontsize{28}{32}\selectfont\bfseries] at (5.3,5.45) {Et après ?};
    \node[text=white, font=\fontsize{18}{22}\selectfont\bfseries, align=center] at (5.3,4.0) {Liste de listes (de listes\\ de listes de ....)};
    \node[text=white, font=\fontsize{18}{22}\selectfont\bfseries, align=center] at (5.3,1.8) {Et si on faisait un jeu de\\ \og Morpion \fg ?};
    \node[fill=cmgold, anchor=south east, font=\footnotesize\bfseries, inner sep=3pt] at (10.2,0.3) {Démo Thonny : brouillon\_morpion.py};
  \end{tikzpicture}
\end{frame}

\begin{frame}{Revenons à notre morpion}

  Démo sur Thonny

  \pause

  \begin{block}{Les leçons à en tirer}
  \begin{itemize}
    \item Une liste de listes se comporte comme une matrice : \mintinline{python}{grille[i][j]} permet d'accéder à la case $(i, j)$ \newline (et la modifier)
    \item Lorsqu'on utilise plusieurs fois une liste, c'est la \alert{même} liste même si on a l'impression de l'avoir copiée
    \item Factoriser le code en fonctions permet de le rendre plus lisible
    \item Ces fonctions peuvent \alert{modifier} la liste donnée en argument
  \end{itemize}
  \end{block}
\end{frame}

\begin{frame}[fragile]{La liste de listes comme matrice}

  \begin{columns}[T]
    \begin{column}{0.62\textwidth}
  \begin{mdframed}[roundcorner=5pt]
\begin{minted}{python}
grille = [
  #   0    1    2
    ["a", "b", "c"], # 0
    ["d", "e", "f"], # 1
    ["g", "h", "i"]  # 2
]
\end{minted}
\end{mdframed}
    \end{column}
    \begin{column}{0.34\textwidth}
      \begin{center}
      \tableau{a,b,c,d,e,f,g,h,i}{\only<2>{\hlrow{1}}\only<3-5>{\hlcell{1}{2}}}
      \end{center}
    \end{column}
  \end{columns}

\begin{block}{Accéder à une ligne et à une case}
  \begin{columns}
    \begin{column}{0.5\textwidth}
  \begin{mdframed}[roundcorner=5pt]
  \begin{onlyenv}<1,4-5>
\begin{minted}[fontsize=\footnotesize]{python}
l = grille[1]
>>> ["d", "e", "f"]
l[2]
>>> "f"
\end{minted}
  \end{onlyenv}
  \begin{onlyenv}<2>
\begin{minted}[fontsize=\footnotesize, highlightlines={1}, highlightcolor=yellow!45]{python}
l = grille[1]
>>> ["d", "e", "f"]
l[2]
>>> "f"
\end{minted}
  \end{onlyenv}
  \begin{onlyenv}<3>
\begin{minted}[fontsize=\footnotesize, highlightlines={3}, highlightcolor=yellow!45]{python}
l = grille[1]
>>> ["d", "e", "f"]
l[2]
>>> "f"
\end{minted}
  \end{onlyenv}
\end{mdframed}
\end{column}
    \begin{column}{0.5\textwidth}
  \begin{mdframed}[roundcorner=5pt]
  \begin{onlyenv}<1-3>
\begin{minted}[fontsize=\footnotesize]{python}
(grille[1])[2]
>>> "f"
grille[1][2]
>>> "f"
\end{minted}
  \end{onlyenv}
  \begin{onlyenv}<4>
\begin{minted}[fontsize=\footnotesize, highlightlines={1}, highlightcolor=yellow!45]{python}
(grille[1])[2]
>>> "f"
grille[1][2]
>>> "f"
\end{minted}
  \end{onlyenv}
  \begin{onlyenv}<5>
\begin{minted}[fontsize=\footnotesize, highlightlines={3}, highlightcolor=yellow!45]{python}
(grille[1])[2]
>>> "f"
grille[1][2]
>>> "f"
\end{minted}
  \end{onlyenv}
    \end{mdframed}
\end{column}
\end{columns}
\end{block}

\end{frame}

\begin{frame}[fragile]{La liste de listes comme matrice}

  \begin{columns}[T]
    \begin{column}{0.62\textwidth}
  \begin{mdframed}[roundcorner=5pt]
\begin{minted}{python}
grille = [
  #   0    1    2
    ["a", "b", "c"], # 0
    ["d", "e", "f"], # 1
    ["g", "h", "i"]  # 2
]
\end{minted}
\end{mdframed}
    \end{column}
    \begin{column}{0.34\textwidth}
      \begin{center}
      \only<1>{\tableau{a,b,c,d,e,f,g,h,i}{\hlrow{1}}}\only<2-3>{\tableau{a,b,c,d,e,j,g,h,i}{\hlcell{1}{2}}}
      \end{center}
    \end{column}
  \end{columns}
\begin{block}{Modifier une case}
  \begin{mdframed}[roundcorner=5pt]
  \begin{onlyenv}<1>
\begin{minted}[fontsize=\footnotesize, highlightlines={1}, highlightcolor=yellow!45]{python}
l = grille[1]
>>> ["d", "e", "f"]
l[2] = "j"
print(grille)
>>> [['a', 'b', 'c'], ['d', 'e', 'j'], ['g', 'h', 'i']]
\end{minted}
  \end{onlyenv}
  \begin{onlyenv}<2>
\begin{minted}[fontsize=\footnotesize, highlightlines={3}, highlightcolor=yellow!45]{python}
l = grille[1]
>>> ["d", "e", "f"]
l[2] = "j"
print(grille)
>>> [['a', 'b', 'c'], ['d', 'e', 'j'], ['g', 'h', 'i']]
\end{minted}
  \end{onlyenv}
  \begin{onlyenv}<3>
\begin{minted}[fontsize=\footnotesize, highlightlines={4}, highlightcolor=yellow!45]{python}
l = grille[1]
>>> ["d", "e", "f"]
l[2] = "j"
print(grille)
>>> [['a', 'b', 'c'], ['d', 'e', 'j'], ['g', 'h', 'i']]
\end{minted}
  \end{onlyenv}
\end{mdframed}
\end{block}
\end{frame}

\begin{frame}[fragile]{La liste de listes comme matrice}

  \begin{columns}[T]
    \begin{column}{0.62\textwidth}
  \begin{mdframed}[roundcorner=5pt]
\begin{minted}{python}
grille = [
  #   0    1    2
    ["a", "b", "c"], # 0
    ["d", "e", "j"], # 1
    ["g", "h", "i"]  # 2
]
\end{minted}
\end{mdframed}
    \end{column}
    \begin{column}{0.34\textwidth}
      \begin{center}
      \only<1>{\tableau{a,b,c,d,e,j,g,h,i}{}}\only<2-3>{\tableau{a,b,c,d,e,j,g,h,k}{\hlcell{2}{2}}}\only<4>{\tableau{a,l,c,d,e,j,g,h,k}{\hlcell{0}{1}}}
      \end{center}
    \end{column}
  \end{columns}
\begin{block}{Modifier une case (suite)}
  \begin{mdframed}[roundcorner=5pt]
  \begin{onlyenv}<1>
\begin{minted}[fontsize=\footnotesize]{python}
(grille[2])[2] = "k"
print(grille)
>>> [['a', 'b', 'c'], ['d', 'e', 'j'], ['g', 'h', 'k']]
grille[0][1] = "l"
>>> [['a', 'l', 'c'], ['d', 'e', 'j'], ['g', 'h', 'k']]
\end{minted}
  \end{onlyenv}
  \begin{onlyenv}<2>
\begin{minted}[fontsize=\footnotesize, highlightlines={1}, highlightcolor=yellow!45]{python}
(grille[2])[2] = "k"
print(grille)
>>> [['a', 'b', 'c'], ['d', 'e', 'j'], ['g', 'h', 'k']]
grille[0][1] = "l"
>>> [['a', 'l', 'c'], ['d', 'e', 'j'], ['g', 'h', 'k']]
\end{minted}
  \end{onlyenv}
  \begin{onlyenv}<3>
\begin{minted}[fontsize=\footnotesize, highlightlines={2}, highlightcolor=yellow!45]{python}
(grille[2])[2] = "k"
print(grille)
>>> [['a', 'b', 'c'], ['d', 'e', 'j'], ['g', 'h', 'k']]
grille[0][1] = "l"
>>> [['a', 'l', 'c'], ['d', 'e', 'j'], ['g', 'h', 'k']]
\end{minted}
  \end{onlyenv}
  \begin{onlyenv}<4>
\begin{minted}[fontsize=\footnotesize, highlightlines={4}, highlightcolor=yellow!45]{python}
(grille[2])[2] = "k"
print(grille)
>>> [['a', 'b', 'c'], ['d', 'e', 'j'], ['g', 'h', 'k']]
grille[0][1] = "l"
>>> [['a', 'l', 'c'], ['d', 'e', 'j'], ['g', 'h', 'k']]
\end{minted}
  \end{onlyenv}
    \end{mdframed}
\end{block}
\end{frame}

\begin{frame}[fragile]{La copie superficielle}

  \begin{mdframed}[roundcorner=5pt]
\begin{minted}[fontsize=\footnotesize]{python}
ligne = ["a", "b", "c"]
grille = [ligne, ligne, ligne]
print(grille)
>>> [["a", "b", "c"], ["a", "b", "c"], ["a", "b", "c"]]
\end{minted}
\end{mdframed}
\begin{block}{Modifier une (des ?) case}
  \begin{mdframed}[roundcorner=5pt]
  \begin{minted}[escapeinside=||, fontsize=\footnotesize]{python}
ligne[1] = "d"
print(grille)
|\pause|>>> [['a', 'd', 'c'], ['a', 'd', 'c'], ['a', 'd', 'c']]
|\pause|grille[0][2] = "e"
|\pause|>>> [['a', 'd', 'e'], ['a', 'd', 'e'], ['a', 'd', 'e']]
    \end{minted}
    \end{mdframed}
\end{block}
La preuve sur \href{https://pythontutor.com/}{Python Tutor}.
\end{frame}

\section{Les fonctions sur les listes}

\begin{frame}[fragile]{Mutabilité, immutabilité}

  \begin{block}{Rappel}
    \vspace{0pt}
  Une fonction \alert{ne peut pas modifier} une variable de type \mintinline{python}{int} donnée en argument :
\vspace{-0.5em}
  \begin{mdframed}[roundcorner=5pt]
  \begin{minted}{python}
def essaie_de_modifier_pour_voir(x):
    x = 5

x = 12
essaie_de_modifier_pour_voir(x)
print(x)
>>> 12
\end{minted}
\end{mdframed}

\vspace{-0.5em}
C'est aussi vrai pour tous les types ``primitifs'' : \newline \mintinline{python}{int}, \mintinline{python}{float}, \mintinline{python}{str}, \mintinline{python}{bool}.

\vspace{-0.5em}
On dit que les types primitifs sont \alert{immutables}.
\vspace{1em}
  \end{block}

  \end{frame}

\begin{frame}[fragile]{Mutabilité, immutabilité}
  Une fonction \alert{peut modifier} la liste donnée en argument :
  \begin{mdframed}[roundcorner=5pt]
  \begin{minted}{python}
def essaie_de_modifier_pour_voir_liste(lst):
    lst[0] = 5

lst = [12]
essaie_de_modifier_pour_voir_liste(lst)
print(lst)
>>> [5]
\end{minted}
\end{mdframed}

On dit que les listes sont \alert{mutables}.
  \end{frame}

  \section{Les itérables}

% Slide refaite en natif (c'était l'image slide_iterables.png)
\begin{frame}{Les itérables}
  \begin{tikzpicture}[x=1cm,y=1cm]
    \useasboundingbox (0,0) rectangle (10.6,6.6);
    \node[align=center, font=\LARGE] at (5.3,5.9) {\textcolor{cmblue}{\bfseries En général :} \textcolor{cmblue}{parcours} de gauche à droite};
    \node[font=\LARGE] at (3.7,4.9) {Élément par élément};
    \node[font=\large\bfseries, text=cmblue] at (3.7,4.1) {OU};
    \node[font=\LARGE] at (3.7,3.3) {Indice par indice};
    \node[fill=cmmint, rounded corners=14pt, anchor=north east, text width=2.5cm, inner sep=8pt, font=\small\bfseries, text=black!70] at (10.5,5.4) {%
      \tiret Listes\\[2pt] \tiret Chaînes\\[2pt] \tiret \textit{Tuples}\\[2pt] \tiret \textit{Ensembles}\\[2pt] \tiret \textit{Fichiers\dots}};
    \node[fill=cmblue, font=\Large\bfseries, text=black, inner sep=4pt] at (5.3,1.7) {SEULEMENT POUR LES LISTES ???};
    \node[fill=cmblue, font=\Large\bfseries, text=black, inner sep=4pt] at (5.3,0.8) {NON : \og ITERABLES \fg};
  \end{tikzpicture}
\end{frame}

\end{document}



%%% Local Variables:
%%% TeX-engine: xetex
%%% TeX-command-extra-options: "-shell-escape"
%%% End:
